\documentclass[aspectratio=169, 14pt]{beamer}
\usepackage{beamer_preamble}

\begin{document}

% ==== Title ===============================================================
\titleslide{Lesson 4-4 \textperiodcentered\ Period 1}{Adding \& Subtracting Rational Expressions + Resistance DOK-3 (Q\#15 Prep)}

% ==== Do Now ==============================================================
\begin{frame}{Do Now --- Critique \& Explain}{Algebra 2 \textperiodcentered\ Lesson 4-4 \textperiodcentered\ Single-Period Quick}
\phasebadges{1}{5}{bridge from 4-3}
\vspace{0.2cm}
\begin{projcard}{Do Now (4-4-savvas-model-discuss-lesson-4-4-launch) \textperiodcentered\ Common denominator conflict}{slidegold}
Teo claims that a common denominator of
\(\dfrac{1}{2} + \dfrac{1}{3} + \dfrac{1}{9}\)
is \(2+3+9=14\). Evander claims it is \(2\cdot3\cdot9=54\).

\vspace{0.2cm}
\begin{enumerate}
  \setlength{\itemsep}{2pt}
  \item[\textbf{A.}] Is either student correct? Explain using the word \textit{common denominator}.
  \item[\textbf{B.}] Find the sum. Explain your method.
  \item[\textbf{C.}] Timothy says: multiply denominators together --- always quickest. Do you agree?
\end{enumerate}
\end{projcard}
\end{frame}

% ==== Launch ==============================================================
\begin{frame}{Launch --- Rules + Example 3 + Example 4 (teacher models)}{Algebra 2 \textperiodcentered\ Lesson 4-4 \textperiodcentered\ Single-Period Quick}
\phasebadges{1--2}{13}{teacher models}
\vspace{0.15cm}
\begin{projcard}{Rules --- Add/Subtract Rational Expressions (keep printed)}{slidegreen}
\small
\begin{enumerate}
  \setlength{\itemsep}{0pt}
  \item LIKE denominators: combine numerators, keep denominator, simplify.
  \item UNLIKE denominators: (a) factor each denominator, (b) find the LCM,
        (c) rewrite each fraction with LCM as denominator, (d) combine numerators, (e) simplify.
  \item LCM = every DISTINCT factor at its highest power across both denominators.
  \item Always check: can the final expression simplify by cancelling a common factor?
\end{enumerate}

\vspace{0.1cm}
\textcolor{slideaccent}{\textbf{COMPRESSED: Teacher models Ex 3 (add unlike) AND Ex 4 (subtract unlike). Watch the minus sign on Ex 4 --- it distributes across EVERY term.}}
\end{projcard}
\end{frame}

% ==== Explore =============================================================
\begin{frame}{Explore --- TPS + DOK-3 Capstone}{Algebra 2 \textperiodcentered\ Lesson 4-4 \textperiodcentered\ Single-Period Quick}
\phasebadges{2--3}{32}{student work}
\small\textcolor{slidegray}{6 items in order. Plan \(\sim\)10 min for Practice \#32 (Resistance). Wed 45-min variant: skip \#16 + \#26.}

\vspace{0.1cm}
\begin{projcard}{Try It 3 \textperiodcentered\ Add unlike denominators}{slideblue}
Find the sum of two rational expressions with unlike denominators. Factor first; build the LCM.
\end{projcard}
\vspace{-0.12cm}
\begin{projcard}{Try It 4 \textperiodcentered\ Subtract unlike denominators}{slideblue}
Simplify. Watch: minus sign distributes across EVERY term of the second numerator.
\end{projcard}
\vspace{-0.12cm}
\begin{projcard}{Practice \#12 \textperiodcentered\ Like denominators}{slideblue}
\textbf{Generalize} --- how is adding rational expressions similar to and different from adding rational numbers?
\end{projcard}
\vspace{-0.12cm}
\begin{projcard}{Practice \#16 \textperiodcentered\ LCM of polynomial denominators}{slideblue}
Find the LCM of \(3x^2+7x+2\) and \(9x+3\). Factor first.
\end{projcard}
\vspace{-0.12cm}
\begin{projcard}{Practice \#26 \textperiodcentered\ Sign-distribution trap (bike ride)}{slideblue}
Ahmed's 15-mile bike trail. Speed first half \(= x+2\); second half \(= x\). Total travel time expression.
\end{projcard}
\vspace{-0.12cm}
\begin{projcard}{Practice \#32 \textperiodcentered\ DOK-3 ELECTRICAL RESISTANCE (\(\bigstar\))}{slideaccent}
Find the ratio of resistance of Circuit A to Circuit B. Setup: \(1/R_{\text{total}} = 1/R_1 + 1/R_2\). Build LCM, combine, form the ratio as a single simplified rational expression.
\end{projcard}
\end{frame}

% ==== Share / Summary =====================================================
\begin{frame}{Share / Summary}{Algebra 2 \textperiodcentered\ Lesson 4-4 \textperiodcentered\ Single-Period Quick}
\phasebadges{2}{5}{DOK-3 reveal + Q\#15 bridge}
\vspace{0.2cm}
\begin{projcard}{Sentence frame \textperiodcentered\ Resistance (\#32)}{slideblue}
``For the parallel circuit, \(1/R_{\text{total}} = 1/R_1 + 1/R_2\). I found a common denominator of \underline{\hspace{1cm}}, added the reciprocals to get \(1/R_{\text{total}} =\) \underline{\hspace{1.5cm}}, then flipped to get \(R_{\text{total}} =\) \underline{\hspace{1cm}}. The ratio \(R_A/R_B\) simplified to \underline{\hspace{1cm}} because \underline{\hspace{2cm}}.''
\end{projcard}

\vspace{0.15cm}
\begin{projcard}{Bridge to Topic 4 LEHS Q\#15}{slideaccent}
The resistance move --- find LCM, combine reciprocals, form a single rational expression --- is the structural reasoning tested on Topic 4 LEHS Q\#15. You have now rehearsed it.
\end{projcard}
\end{frame}

% ==== Exit Ticket =========================================================
\begin{frame}{Exit Ticket --- Summary Recap}{Algebra 2 \textperiodcentered\ Lesson 4-4 \textperiodcentered\ Single-Period Quick}
\phasebadges{1--2}{3}{not graded}
\vspace{0.2cm}
\begin{projcard}{Complete on your exit ticket}{slidegold}
\begin{enumerate}
  \setlength{\itemsep}{0.25cm}
  \item To add rational expressions with unlike denominators, I must first find the \underline{\hspace{4cm}} of the denominators.
  \item Practice \#32 (resistance) showed me that reciprocal-sum problems (like \(1/R_1 + 1/R_2\)) use the same algebraic steps as adding rational expressions because \underline{\hspace{4cm}}.
  \item One biggest thing I learned today: \underline{\hspace{6cm}}
\end{enumerate}
\end{projcard}
\end{frame}

\end{document}
